% defined-in: zorich (page 319), zorich (page 543), zorich (page 319), zorich (page 319)
% === SEMANTIC MAPPING ===
% 1. Variables & Types: 
%    - a, b : \RealNumbers (a < b)
%    - f : [a, b] \to \RealNumbers
%    - I : \RealNumbers
%    - n : \NaturalNumbers
%    - P : \text{Partition}([a, b]) \equiv \{x : \text{Fin}(n+1) \to \RealNumbers \mid x(0) = a \land x(n) = b \land \forall i < n, x(i) < x(i+1)\}
%    - \xi : \text{MarkedPoints}(P) \equiv \{\xi : \text{Fin}(n) \to \RealNumbers \mid \forall i < n, x(i) \leq \xi(i) \leq x(i+1)\}
%    - \lambda : \text{Partition} \to \RealNumbers \equiv \lambda(P) = \max_{i < n} (x(i+1) - x(i))
%    - \sigma : \text{Function} \times \text{Partition} \times \text{MarkedPoints} \to \RealNumbers \equiv \sum_{i=0}^{n-1} f(\xi(i)) \cdot (x(i+1) - x(i))
% 2. Data Structures: 
%    - Partition represented as a strictly increasing function from a finite index set to \RealNumbers.
%    - Marked points represented as a function from the index set of sub-intervals to \RealNumbers.
% 3. Implicit Bounds: 
%    - n : \NaturalNumbers is the number of sub-intervals.
%    - i : \text{Fin}(n) is the index of sub-intervals.
% ========================

% entity-id: def-riemann-integral
% entity-type: def
% macro: \RiemannIntegral
% args: 1
% notation: #1
\hypertarget{def-riemann-integral}{}
\begin{definition}[def-riemann-integral]
\[
\TerForall \Function{f} \colon \ClosedInterval{a,b} \TermTo \RealNumbers
\TerForall I \colon \RealNumbers
\mathrm{IsRiemannIntegral}(\Function{f}, I) \TermDefinedAs
\TerForall \varepsilon \colon \RealNumbers \left( \varepsilon > 0 \TermImplication 
\TermExists \delta \colon \RealNumbers \left( \delta > 0 \TermConjunction 
\TerForall n \colon \NaturalNumbers \TerForall P \TermIn \Partition{[a,b]} \TerForall \xi \TermIn \Partition{P} 
\left( \Partition{\lambda}(P) < \delta \TermImplication 
\AbsAbstract{ I - \sum_{i=0}^{n-1} f(\xi(i)) \cdot (x(i+1) - x(i)) } < \varepsilon \right) \right) \right)
\]
\end{definition}